\subsection{§III.2.1 (Déf.~2) --- Transitivité des segments}
\textbf{Énoncé (Bourbaki).} Soit $E$ un ensemble ordonné par une relation $\le$. Si $S$ est un segment de $E$ et $T$ un segment de $S$, alors $T$ est un segment de $E$. Un segment $A\subset E$ est une partie close vers le bas : $A\subset E$ et $(\forall x)(\forall y)\bigl(((x\in A \wedge y\in E)\wedge y\le x)\Rightarrow y\in A\bigr)$.

\textbf{Formalisé.} Cible : $\seg(T,R,E)$, c.-à-d. $T\subset E \;\wedge\; (\forall x)(\forall y)\bigl(((x\in T\wedge y\in E)\wedge y\le x)\Rightarrow y\in T\bigr)$, avec la convention de relation $y\le x := R(y,x)$. \\
Module : \texttt{bourbaki/ordre/iii\_2\_bon\_ordre/bon\_ordre\_segments/ensembles\_segment\_transitif.py}.

\textbf{Preuve (\Nstar).} On \texttt{assume} les deux énoncés de segment, puis \texttt{conjonction\_elim\_gauche}/\texttt{droite} en sépare les composantes : inclusions $S\subset E$, $T\subset S$ et clôtures $\mathrm{clos}_S$, $\mathrm{clos}_T$. Composante inclusion : la tactique \texttt{inclusion\_transitive} en $(T,S,E)$ donne $((T\subset S)\wedge(S\subset E))\Rightarrow(T\subset E)$, fermée par \texttt{conjonction\_intro} + \texttt{modus\_ponens} en $T\subset E$. Composante clôture : sous \texttt{assume} de la prémisse $((x\in T\wedge y\in E)\wedge y\le x)$, on chaîne $x\in T \Rightarrow x\in S$ (\texttt{instancie} $T\subset S$ + \texttt{modus\_ponens}), puis $\mathrm{clos}_S$ instancié en $(x,y)$ donne $y\in S$ et $\mathrm{clos}_T$ donne $y\in T$ (motif \texttt{instancie}$\times 2$ puis \texttt{modus\_ponens} sur la prémisse reconstruite par \texttt{conjonction\_intro}). \texttt{loi\_deduction} décharge la prémisse, \texttt{generalisation} en $y$ puis $x$ ferme le corps ; \texttt{conjonction\_intro} recolle inclusion et clôture en $\seg(T,R,E)$.

\textbf{Statut.} CLOS sous hypothèses honnêtes $\{\,\seg(S,R,E),\ \seg(T,R,S)\,\}$ ; \texttt{theorie\_ensembles()==22}. Séquent exact (les deux hypothèses ne sont jamais déchargées, aucune parasite) ; la forme implication fermée $\vdash (\seg(S,R,E)\wedge\seg(T,R,S))\Rightarrow\seg(T,R,E)$ s'en déduit en une ligne par \texttt{loi\_deduction}.

\subsection{§III.2.1 (Prop.~2, préliminaire) --- $x\notin \seg_x$ et témoin de $\seg_x \subsetneq \seg_y$}
\textbf{Énoncé (Bourbaki).} L'application $x\mapsto \seg_x$ (segment initial strict d'extrémité $x$) est strictement croissante pour $\subset$. La monotonie large $\seg_x\subset\seg_y$ dès que $x\le y$ étant acquise ailleurs, la strictité $\seg_x\subsetneq\seg_y$ pour $x\neq y$ requiert un témoin d'écart : un élément de $\seg_y$ hors de $\seg_x$. Ce témoin est $x$ lui-même. Le segment strict est $\seg_x := \{u\in E \mid u\le x \wedge u\neq x\}$, caractérisé par l'axiome de segment $u\in\seg_x \iff ((u\in E \wedge u\le x)\wedge u\neq x)$.

\textbf{Formalisé.} Deux cibles. (a) $\neg(x\in\seg(R,E,x))$, INCONDITIONNEL. (b) $\bigl(x\in\seg(R,E,y)\bigr)\wedge\neg\bigl(x\in\seg(R,E,x)\bigr)$, le témoin d'écart. \\
Module : \texttt{bourbaki/ordre/iii\_2\_bon\_ordre/bon\_ordre\_segments/ensembles\_segment\_strict\_propre.py}.

\textbf{Preuve (\Nstar).} (a) \texttt{element\_hors\_de\_son\_segment} : on \texttt{assume} $x\in\seg_x$ ; \texttt{equivalence\_avant} sur l'axiome de segment (\texttt{membre\_segment} en $(x,x)$) + \texttt{modus\_ponens} projette par \texttt{conjonction\_elim\_droite} la composante $x\neq x$. Confrontée à $x=x$ (\texttt{N.reflexivite}), un ex falso local (via \texttt{N.s2}) vise directement $\neg(x\in\seg_x)$ ; \texttt{loi\_deduction} décharge l'hypothèse et un \texttt{N.s1} de réfutation ($P\Rightarrow\neg P \vdash \neg P$) clôt. (b) \texttt{seg\_strict\_propre} : la composante $x\in\seg_y$ s'obtient en \texttt{assume}-ant $x\in E$, $x\le y$, $x\neq y$, reconstruits par \texttt{conjonction\_intro} en le corps de l'axiome, puis \texttt{equivalence\_arriere} (\texttt{membre\_segment} sens arrière) + \texttt{modus\_ponens} ; la composante $x\notin\seg_x$ est le lemme (a) ; \texttt{conjonction\_intro} assemble le témoin.

\textbf{Statut.} (a) CLOS (0 hypothèse) : \texttt{est\_clos} vérifié par assertion, tout dérive de l'axiome de segment et de la réflexivité de l'égalité (déjà présents). (b) CLOS sous hypothèses honnêtes $\{\,x\in E,\ R\{x,y\},\ x\neq y\,\}$ : exactement les antécédents load-bearing pour $x\in\seg_y$ ; \texttt{est\_relation\_ordre}$(R)$ n'est pas consommé, donc absent. \texttt{theorie\_ensembles()==22} ; aucune conclusion n'est l'une de ses hypothèses.

\subsection{§III.2.1 --- Le segment du plus petit élément est vide}
\textbf{Énoncé (Bourbaki).} Soit $E$ bien ordonné par $\le$ et $\alpha = \min E$ son plus petit élément. Alors le segment initial $\seg_\alpha = {]\!\leftarrow,\alpha[}$ d'extrémité $\alpha$ est vide. En effet, $u\in\seg_\alpha$ exigerait $u\le\alpha$ et $u\neq\alpha$ ; or $\alpha$ minore $E$, donc $\alpha\le u$, et l'antisymétrie force $u=\alpha$, contradiction.

\textbf{Formalisé.} Cible : $\seg(R,E,\alpha)=\emptyset$, avec $R$ porté comme graphe ($R\{a,b\} := (a,b)\in R$). \\
Module : \texttt{bourbaki/ordre/iii\_2\_bon\_ordre/bon\_ordre\_segments/ensembles\_segment\_minimum.py}.

\textbf{Preuve (\Nstar).} On \texttt{assume} \texttt{est\_bien\_ordonne}$(R,E)$ et \texttt{est\_plus\_petit\_element}$(R,E,\alpha)$. Des \texttt{conjonction\_elim} en cascade extraient l'antisymétrie $\mathrm{ordre\_antisymetrique}(R)$ du premier et le minorant $(\forall x)(x\in E\Rightarrow R\{\alpha,x\})$ du second. Pour montrer $(\forall u)\neg(u\in\seg_\alpha)$ : sous \texttt{assume} $u\in\seg_\alpha$, \texttt{equivalence\_avant} (\texttt{membre\_segment}) + \texttt{modus\_ponens} projettent $u\in E$, $R\{u,\alpha\}$, $u\neq\alpha$. Le minorant \texttt{instancie} en $u$ donne $R\{\alpha,u\}$ ; l'antisymétrie \texttt{instancie}$\times 2$ en $(u,\alpha)$, appliquée par \texttt{conjonction\_intro} + \texttt{modus\_ponens} à $R\{u,\alpha\}\wedge R\{\alpha,u\}$, force $u=\alpha$ : ex falso local avec $u\neq\alpha$, \texttt{loi\_deduction} et réfutation donnent $\neg(u\in\seg_\alpha)$. Enfin $\seg_\alpha=\emptyset$ par extensionnalité \textbf{A1} : $\seg_\alpha\subset\emptyset$ (\texttt{N.s2} depuis $\neg(u\in\seg_\alpha)$, \texttt{generalisation}) et $\emptyset\subset\seg_\alpha$ (ex falso depuis \texttt{AXIOME\_VIDE}), recombinés par \texttt{conjonction\_intro} + \texttt{modus\_ponens} sur \textbf{A1} instancié en $(\seg_\alpha,\emptyset)$.

\textbf{Statut.} CLOS sous hypothèses honnêtes $\{\,$\texttt{est\_bien\_ordonne}$(R,E)$, \texttt{est\_plus\_petit\_element}$(R,E,\alpha)\,\}$ ; \texttt{theorie\_ensembles()==22}. Séquent exact (les deux hypothèses sont réellement consommées --- antisymétrie et minorant ---, aucune parasite, la conclusion n'est aucune hypothèse). Rien postulé : tout dérive de l'axiome de segment, de \textbf{A1} et d'\texttt{AXIOME\_VIDE}, déjà présents.

Les résultats nommés \texttt{segment\_extremite\_strictement\_croissant} (conjonction des deux briques ci-dessus) et la ré-exposition de la monotonie large complètent la Proposition~2.
